\begin{itemize}
    \item Medical Linear Accelerator
    \item 3D Scanner (RFA) with Water Reservoir
    \item Two Ionization Chamber (Field and reference Ionization chamber)
    \item Electrometer and Connecting cables.
    \item Thermometer and Barometer
    \item Leveling tool (Spirit level)
\end{itemize}
\subsection{Medical Linear Accelerator}
A medical linear accelerator (linac) is a device that uses high-frequency electromagnetic waves to accelerate charged particles, such as electrons, to high energies. These accelerated particles are then used to produce high-energy X-rays or electron beams for radiation therapy. The linac consists of several components, including an electron gun, a waveguide, and a treatment head. The electron gun generates electrons, which are then accelerated through the waveguide by the electromagnetic waves. The treatment head contains various components, such as a target for X-ray production, collimators to shape the beam, and a multileaf collimator for further beam shaping. The linac is a crucial tool in modern radiation therapy, allowing for precise targeting of tumors while minimizing damage to surrounding healthy tissue.
\subsection{3D Scanner (RFA) with Water Reservoir}
A 3D Scanner with a water reservoir, often referred to as a Radiation Field Analyzer (RFA), is a device used in radiation therapy to measure and analyze the characteristics of radiation beams. The water reservoir is filled with water, which serves as a medium for simulating human tissue. The RFA is equipped with a scanning mechanism that allows it to move through the water, measuring the dose distribution of the radiation beam at various depths and positions. This information is crucial for treatment planning and quality assurance in radiation therapy, as it helps ensure that the prescribed dose is delivered accurately to the target area while minimizing exposure to surrounding healthy tissues. RFA is a large, usually 50 x 50 x 50cm3 water phantom used to perform radiation dosimetry in
a clinical setup. The RFA can hold 166 L of water, and it has a railing on which the chamber can be
fixed using holders and moved using a control setup from outside. 3 screws are provided at the base
of the RFA to level it. The whole mechanical system is such that a very high accuracy of chamber
positioning (0.4 mm) can be achieved (as the ion chambers can move in 3 orthogonal directions).The RFA can provide detailed data on the beam's intensity, shape, and uniformity, which are essential for optimizing treatment outcomes.
\subsection{Ionization Chambers}
Ionization chambers are devices used to measure ionizing radiation. They consist of a gas-filled chamber with two electrodes. When ionizing radiation passes through the chamber, it ionizes the gas, creating positive ions and free electrons. The electrodes collect these charges, producing an electrical signal that is proportional to the amount of radiation. In radiation therapy, two types of ionization chambers are commonly used: the field ionization chamber and the reference ionization chamber. The field ionization chamber is used to measure the dose distribution of the radiation beam as it scans through the water in the RFA. The reference ionization chamber is typically fixed in a corner of the water tank and monitors the beam output in real-time. The signal from the field ionization chamber is normalized against the reference chamber's signal to remove errors caused by fluctuations in the machine's dose rate, ensuring accurate dose measurements for treatment planning and quality assurance.
\subsection{Electrometer and Connecting Cables}
An electrometer is a sensitive electronic device used to measure electric charge or current. In the context of radiation therapy, electrometers are used to measure the electrical signals produced by ionization chambers when they detect ionizing radiation. The electrometer amplifies and converts these signals into readable values, which can then be used to calculate the dose of radiation being delivered. Connecting cables are essential for linking the ionization chambers to the electrometer, allowing for the transmission of electrical signals. These cables must be of high quality to ensure accurate signal transmission without interference or loss, which is crucial for obtaining precise measurements in radiation therapy.
\subsection{Thermometer and Barometer}
A thermometer is an instrument used to measure temperature, while a barometer is used to measure atmospheric pressure. In radiation therapy, both of these instruments are important for ensuring accurate dose measurements. The temperature and pressure of the environment can affect the density of the air in the ionization chambers, which in turn can influence the readings obtained from these chambers. By monitoring and accounting for changes in temperature and pressure, clinicians can apply necessary corrections to the dose measurements, ensuring that the prescribed dose is delivered accurately to the patient. This is particularly important in radiation therapy, where precise dosing is critical for effective treatment and minimizing side effects.
\subsection{Leveling Tool (Spirit Level)}A leveling tool, such as a spirit level, is used to ensure that the equipment and setup in radiation therapy are properly aligned and level. Proper leveling is crucial for accurate dose delivery, as any tilt or misalignment can lead to uneven dose distribution and potentially harm healthy tissues while underdosing the target area. The spirit level typically consists of a sealed tube containing a liquid and an air bubble. When the tool is placed on a surface, the position of the bubble indicates whether the surface is level. In radiation therapy, the spirit level can be used to check the alignment of the treatment couch, the water reservoir, and other equipment to ensure that everything is properly positioned for accurate dose delivery. This helps to maintain the precision and effectiveness of the treatment, ultimately improving patient outcomes.